\appendix

\section{Proof Details}

\subsection{Gaussian initialization and fixed-span preservation}\label{app:step-001}

\begin{lemma}[Admissible dimensions and Gaussian full rank]\label{lem:step-001-full-rank}
Under Assumptions~\ref{assump:dimension}, \ref{assump:rank_window}, and
\ref{assump:joint_initialization},
\[
 \frac{k}{n}\le\frac18<1,
\]
and
\[
 \mathbb P\!\left(
 \operatorname{rank}G_x^{\mathrm{cALS}}=k
 \text{ and }
 \operatorname{rank}G_x^{\mathrm{cGD}}=k
 \right)=1.
\]
\end{lemma}

\begin{proof}
Assumptions~\ref{assump:dimension} and \ref{assump:rank_window} give
\[
 0<k\le r^{5/4}\le\frac n8,
\]
so $k/n\le1/8$ and $k<n$.

We first recall the elementary polynomial fact needed below.  The zero set of
a nonzero real polynomial has Lebesgue measure zero.  In one variable this
follows because a nonzero polynomial has finitely many roots.  Inductively,
write a polynomial in $d$ variables as a polynomial in its last variable,
with coefficient polynomials in the first $d-1$ variables.  At least one
coefficient is nonzero.  The common zero set of the nonzero coefficients has
measure zero by induction; outside it, the last-variable polynomial is
nonzero and has finitely many roots.  Fubini's theorem proves the assertion in
dimension $d$.

Fix $M\in\{\mathrm{cALS},\mathrm{cGD}\}$ and let $B_M$ be the leading
$k\times k$ block of $G_x^M$, which exists because $k<n$.  The polynomial
$B\mapsto\det B$ is not identically zero, since it equals one at $I_k$.
The entries of $B_M$ have a joint density, and therefore
\[
 \mathbb P(\det B_M=0)=0.
\]
Whenever $B_M$ is nonsingular, the $k$ columns of $G_x^M$ are linearly
independent.  Thus each matrix has column rank $k$ almost surely.  A finite
union of the two null complements proves the simultaneous statement;
independence is not needed for this last intersection.
\end{proof}

\begin{lemma}[Haar initialization spans and independence]\label{lem:step-001-haar-independence}
Under Assumptions~\ref{assump:dimension}, \ref{assump:rank_window}, and
\ref{assump:joint_initialization}, on the probability-one event of
Lemma~\ref{lem:step-001-full-rank}, each
\[
 \mathcal S_M=\operatorname{range}(G_x^M),
 \qquad M\in\{\mathrm{cALS},\mathrm{cGD}\},
\]
is Haar-uniform on the Grassmannian $\operatorname{Gr}(k,n)$.  The random
elements $\mathcal S_{\mathrm{cALS}}$, $\mathcal S_{\mathrm{cGD}}$, and $T$
are mutually independent.  A regular conditional law of the pair of spans
given $T=T_0$ may be chosen, for every $T_0$, as
\[
 \operatorname{Haar}_{k,n}\otimes\operatorname{Haar}_{k,n}.
\]
The corresponding subspace projectors, tensor subspaces, and tensor
projectors inherit these independence properties.
\end{lemma}

\begin{proof}
Fix a method $M$ and $U\in O(n)$.  If $g_i$ is a column of $G_x^M$, then
\[
 Ug_i\sim\mathcal N\!\left(0,U(n^{-1}I_n)U^{\mathsf T}\right)
 =\mathcal N(0,n^{-1}I_n).
\]
The columns remain independent, and hence
\[
 UG_x^M\stackrel{d}=G_x^M.
\]
On the full-rank event,
\[
 \operatorname{range}(UG_x^M)=U\operatorname{range}(G_x^M)=U\mathcal S_M,
\]
so the law of $\mathcal S_M$ is invariant under every orthogonal
transformation.

For completeness, this invariant law is unique.  Let $\lambda$ be normalized
Haar measure on the compact group $O(n)$ and let $\mathcal S_0$ be a fixed
$k$-plane.  The pushforward of $\lambda$ under $U\mapsto U\mathcal S_0$ is
invariant.  If $\mu$ is any other invariant probability on
$\operatorname{Gr}(k,n)$ and $h$ is bounded Borel, then
\[
 \int h(\mathcal S)\,d\mu(\mathcal S)
 =\int\!\int h(U\mathcal S)\,d\lambda(U)\,d\mu(\mathcal S).
\]
For every $\mathcal S=V\mathcal S_0$, right invariance of $\lambda$ makes
the inner integral equal to
$\int h(U\mathcal S_0)\,d\lambda(U)$, independently of $\mathcal S$.
Thus $\mu$ equals the indicated pushforward.  It follows that
$\mathcal S_M$ is Haar-uniform.

Let $\Xi$ denote the complete collection of smoothing variables used to form
$T$.  Assumption~\ref{assump:joint_initialization} makes
\[
 G_x^{\mathrm{cALS}},\qquad G_x^{\mathrm{cGD}},\qquad\Xi
\]
mutually independent.  The spans, represented by their measurable
projectors, are measurable functions of the first two matrices, and $T$ is a
polynomial, hence measurable, function of $\Xi$.  Measurable images preserve
independence, so the two spans and $T$ are mutually independent.

Let $\mu=\operatorname{Haar}_{k,n}\otimes\operatorname{Haar}_{k,n}$.
For Borel sets $A\subseteq\operatorname{Gr}(k,n)^2$ and
$B\subseteq\mathbb R^{n\times n\times n}$,
\[
 \begin{aligned}
 &\mathbb P\!\left(
 (\mathcal S_{\mathrm{cALS}},\mathcal S_{\mathrm{cGD}})\in A,
 T\in B\right)\\
 &\qquad=\mu(A)\mathbb P(T\in B)
 =\int_B\mu(A)\,\mathbb P_T(dT_0).
 \end{aligned}
\]
Therefore the constant kernel $T_0\mapsto\mu$ is a regular conditional law
for every $T_0$.  The projector and tensor-subspace assertions follow because
they are measurable functions of the spans.
\end{proof}

\begin{lemma}[First-factor membership lifts to tensor membership]\label{lem:step-001-tensor-lift}
Let $Q\in\mathbb R^{n\times k}$ have orthonormal columns, set
$\mathcal S=\operatorname{range}(Q)$, and define
\[
 \mathcal H=\mathcal S\otimes\mathbb R^n\otimes\mathbb R^n,
 \qquad
 P_{\mathcal H}=QQ^{\mathsf T}\otimes I_n\otimes I_n.
\]
If $D\in\mathbb R^{k\times k}$, $X=QD$, and
$Y,Z\in\mathbb R^{n\times k}$, then
\[
 S(X,Y,Z)\in\mathcal H,
 \qquad
 P_{\mathcal H}S(X,Y,Z)=S(X,Y,Z),
 \qquad
 (I-P_{\mathcal H})S(X,Y,Z)=0.
\]
\end{lemma}

\begin{proof}
Write $D=[d_1,\ldots,d_k]$.  Then $x_i=Qd_i\in\mathcal S$, and hence
\[
 x_i\otimes y_i\otimes z_i
 \in\mathcal S\otimes\mathbb R^n\otimes\mathbb R^n=\mathcal H
\]
for every $i$.  Their sum lies in $\mathcal H$.  Since $QQ^{\mathsf T}$ is
the orthogonal projector onto $\mathcal S$, the tensor-product projector
fixes every summand and therefore fixes the sum.  The complementary-projector
identity is equivalent.
\end{proof}

\begin{proposition}[Exact fixed-span preservation]\label{prop:step-001-fixed-span}
Under Assumptions~\ref{assump:dimension} and \ref{assump:rank_window}, and
under Assumption~\ref{assump:joint_initialization}, consider the
probability-one event in Lemma~\ref{lem:step-001-full-rank}.  Every defined
iterate of either constrained method has a matrix
$D_t^M\in\mathbb R^{k\times k}$ for which
\[
 X_t^M=Q_MD_t^M.
\]
Consequently, for every such index,
\[
 S_t^M\in\mathcal H_M,
 \qquad
 P_{\mathcal H_M}S_t^M=S_t^M,
 \qquad
 (I-P_{\mathcal H_M})S_t^M=0.
\]
\end{proposition}

\begin{proof}
On the simultaneous full-rank event, every $Q_M$ has $k$ orthonormal columns
spanning $\mathcal S_M$.

For constrained ALS, initialization gives
\[
 X_0^{\mathrm{cALS}}
 =Q_{\mathrm{cALS}}Q_{\mathrm{cALS}}^{\mathsf T}
 G_x^{\mathrm{cALS}}
 =Q_{\mathrm{cALS}}D_0^{\mathrm{cALS}},
 \qquad
 D_0^{\mathrm{cALS}}
 =Q_{\mathrm{cALS}}^{\mathsf T}G_x^{\mathrm{cALS}}.
\]
Whenever the next sweep is evaluated, its constrained block has the exact
form
\[
 X_{t+1}^{\mathrm{cALS}}=Q_{\mathrm{cALS}}D_{t+1}^{\mathrm{cALS}},
\]
where
\[
 D_{t+1}^{\mathrm{cALS}}
 =Q_{\mathrm{cALS}}^{\mathsf T}T_{(1)}K_t^x
 \bigl((K_t^x)^{\mathsf T}K_t^x\bigr)^\dagger.
\]
Every finite real matrix has a Moore--Penrose pseudoinverse: an SVD defines
it by reciprocating the nonzero singular values.  Thus this factorization
remains meaningful when the Gram matrix is singular or zero.  The later
$Y$- and $Z$-updates do not alter $X_{t+1}$.

For constrained GD,
\[
 Q_{\mathrm{cGD}}C_0
 =Q_{\mathrm{cGD}}Q_{\mathrm{cGD}}^{\mathsf T}G_x^{\mathrm{cGD}}
 =G_x^{\mathrm{cGD}}=X_0^{\mathrm{cGD}},
\]
and at every later defined index the method defines
$X_t^{\mathrm{cGD}}=Q_{\mathrm{cGD}}C_t$.  Therefore one may take
$D_t^{\mathrm{cGD}}=C_t$.  Lemma~\ref{lem:step-001-tensor-lift} now gives
the three tensor-membership identities for both methods.

This argument is deliberately scoped to defined indices.  It does not assume
termination of future line searches; that all-index existence is proved
separately below.
\end{proof}

\begin{proposition}[Initialization and zero-leakage certificate]\label{prop:step-001-certificate}
Under Assumptions~\ref{assump:dimension}, \ref{assump:rank_window}, and
\ref{assump:joint_initialization}, the two initialization spans are mutually
independent Haar $k$-planes independent of $T$, conditionally retain the
constant product-Haar law given every $T=T_0$, and every defined represented
tensor of either constrained method has zero leakage from its fixed tensor
subspace.
\end{proposition}

\begin{proof}
Lemma~\ref{lem:step-001-full-rank} gives simultaneous full rank with
probability one and the exact arithmetic $k/n\le1/8<1$.
Lemma~\ref{lem:step-001-haar-independence} gives the Haar and conditional
independence statements.  Proposition~\ref{prop:step-001-fixed-span} gives
zero leakage at every defined index.

All boundary regimes are included.  Equality $k/n=1/8$ still gives $k<n$.
The measurable orthonormalization may choose arbitrary signs or orientations,
because only the span and projector are used.  Singular or zero constrained
ALS designs do not affect the left factor $Q_{\mathrm{cALS}}$.  No rank or
size condition on $Y_t,Z_t$ is required by
Lemma~\ref{lem:step-001-tensor-lift}.  At a currently undefined future GD
index no assertion is made, while each existing index has the exact
coefficient representation.  The conclusions do not depend on the value of
$T$ and remain exact at $T=0$ and in exact or zero-smoothing specializations.
Finally, the mechanism is active at initialization and at the first defined
transition: $X_0^M=Q_MD_0^M$, the first constrained ALS update is
$Q_MD_1^M$, and the first constrained GD update changes only its coefficient
matrix.
\end{proof}

\subsection{Conditional Haar projection margin}\label{app:step-002}

\begin{lemma}[Conditional Haar projector isotropy and tensor energy]\label{lem:step-002-projector-energy}
Under Assumptions~\ref{assump:dimension}, \ref{assump:rank_window}, and
\ref{assump:joint_initialization}, and Lemma~\ref{lem:step-001-haar-independence},
fix $M\in\{\mathrm{cALS},\mathrm{cGD}\}$ and
$T_0\in\mathbb R^{n\times n\times n}$.  Under the regular conditional law
given $T=T_0$,
\[
 \mathbb E[P_{\mathcal S_M}\mid T=T_0]=\frac{k}{n}I_n
\]
and
\[
 \mathbb E\!\left[\|P_{\mathcal H_M}T\|_F^2\mid T=T_0\right]
 =\frac{k}{n}\|T_0\|_F^2.
\]
Consequently,
\[
 \mathbb E\!\left[\|P_{\mathcal H_M}T\|_F^2\mid T\right]
 =\frac{k}{n}\|T\|_F^2.
\]
\end{lemma}

\begin{proof}
Let
\[
 \overline P_M=\mathbb E[P_{\mathcal S_M}\mid T=T_0].
\]
Haar invariance gives
$U\overline P_MU^{\mathsf T}=\overline P_M$ for every $U\in O(n)$.
Conjugation by coordinate sign flips makes every off-diagonal entry zero, and
conjugation by coordinate permutations makes all diagonal entries equal.
Thus $\overline P_M=\lambda I_n$.  Since every realization is a rank-$k$
orthogonal projector,
\[
 n\lambda=\operatorname{tr}(\overline P_M)
 =\mathbb E[\operatorname{tr}(P_{\mathcal S_M})\mid T=T_0]=k,
\]
which proves the first identity.

The tensor projector acts in the first mode, so
\[
 (P_{\mathcal H_M}T_0)_{(1)}=P_{\mathcal S_M}(T_0)_{(1)}.
\]
Using symmetry and idempotence,
\[
 \begin{aligned}
 \|P_{\mathcal H_M}T_0\|_F^2
 &=\|P_{\mathcal S_M}(T_0)_{(1)}\|_F^2\\
 &=\operatorname{tr}\!\left(
 P_{\mathcal S_M}(T_0)_{(1)}(T_0)_{(1)}^{\mathsf T}\right).
 \end{aligned}
\]
Taking conditional expectation yields
\[
 \begin{aligned}
 \mathbb E[\|P_{\mathcal H_M}T\|_F^2\mid T=T_0]
 &=\operatorname{tr}\!\left(
 \frac{k}{n}I_n(T_0)_{(1)}(T_0)_{(1)}^{\mathsf T}\right)\\
 &=\frac{k}{n}\|(T_0)_{(1)}\|_F^2
 =\frac{k}{n}\|T_0\|_F^2.
 \end{aligned}
\]
The conditional kernel is valid for every $T_0$, giving the final identity.
\end{proof}

\begin{lemma}[Zero-safe per-method projection event]\label{lem:step-002-markov-event}
Under Assumptions~\ref{assump:dimension}, \ref{assump:rank_window}, and
\ref{assump:joint_initialization}, Lemma~\ref{lem:step-001-haar-independence},
and Lemma~\ref{lem:step-002-projector-energy}, define
\[
 E_M=\left\{
 \|P_{\mathcal H_M}T\|_F^2
 \le2\frac{k}{n}\|T\|_F^2
 \right\}.
\]
Then, for every $T_0$,
\[
 \mathbb P(E_M\mid T=T_0)\ge\frac12.
\]
If $T_0=0$, this conditional probability is one.
\end{lemma}

\begin{proof}
Suppose first that $T_0\ne0$ and set
$X_M=\|P_{\mathcal H_M}T_0\|_F^2\ge0$.  The threshold
$2(k/n)\|T_0\|_F^2$ is positive and
Lemma~\ref{lem:step-002-projector-energy} gives
\[
 \mathbb E[X_M\mid T=T_0]=\frac{k}{n}\|T_0\|_F^2.
\]
Markov's inequality, which follows from
$a\mathbf 1_{\{X_M\ge a\}}\le X_M$ for $a>0$, gives
\[
 \begin{aligned}
 \mathbb P(E_M^c\mid T=T_0)
 &\le\mathbb P\!\left(
 X_M\ge2\frac{k}{n}\|T_0\|_F^2\,\middle|\,T=T_0\right)\\
 &\le\frac{(k/n)\|T_0\|_F^2}
 {2(k/n)\|T_0\|_F^2}=\frac12.
 \end{aligned}
\]
If $T_0=0$, both sides of the inequality defining $E_M$ are zero for every
span, so $E_M$ is sure and no division is used.
\end{proof}

\begin{lemma}[Exact orthogonal-complement margin]\label{lem:step-002-residual-margin}
Under Assumptions~\ref{assump:dimension} and \ref{assump:rank_window}, on
$E_M$ one has
\[
 \|(I-P_{\mathcal H_M})T\|_F^2
 \ge\left(1-2\frac{k}{n}\right)\|T\|_F^2
 \ge\frac34\|T\|_F^2.
\]
The assertion includes $T=0$ and the boundary $k/n=1/8$.
\end{lemma}

\begin{proof}
Orthogonality of $P_{\mathcal H_M}$ gives
\[
 \|T\|_F^2
 =\|P_{\mathcal H_M}T\|_F^2
 +\|(I-P_{\mathcal H_M})T\|_F^2.
\]
On $E_M$,
\[
 \|(I-P_{\mathcal H_M})T\|_F^2
 \ge\left(1-2\frac{k}{n}\right)\|T\|_F^2.
\]
The assumptions give
\[
 \frac{k}{n}\le\frac{r^{5/4}}{8r^{5/4}}=\frac18,
 \qquad
 1-2\frac{k}{n}\ge1-2\left(\frac18\right)=\frac34.
\]
All quantities vanish when $T=0$, and equality at $k/n=1/8$ is allowed.
\end{proof}

\begin{proposition}[Joint conditional and unconditional projection event]\label{prop:step-002-joint-event}
Under Assumptions~\ref{assump:dimension}, \ref{assump:rank_window}, and
\ref{assump:joint_initialization}, Lemma~\ref{lem:step-001-haar-independence},
Lemma~\ref{lem:step-002-markov-event}, and
Lemma~\ref{lem:step-002-residual-margin}, the events
$E_{\mathrm{cALS}}$ and $E_{\mathrm{cGD}}$ are conditionally independent
given $T$.  For every $T_0$,
\[
 \begin{aligned}
 &\mathbb P(E_{\mathrm{cALS}}\cap E_{\mathrm{cGD}}\mid T=T_0)\\
 &\qquad=\mathbb P(E_{\mathrm{cALS}}\mid T=T_0)
 \mathbb P(E_{\mathrm{cGD}}\mid T=T_0)\ge\frac14.
 \end{aligned}
\]
Moreover,
\[
 \mathbb P(E_{\mathrm{cALS}}\cap E_{\mathrm{cGD}})\ge\frac14,
\]
and on this joint event, simultaneously for both methods,
\[
 \|(I-P_{\mathcal H_M})T\|_F^2
 \ge\left(1-2\frac{k}{n}\right)\|T\|_F^2
 \ge\frac34\|T\|_F^2.
\]
\end{proposition}

\begin{proof}
Fix $T_0$.  The event $E_M$ is then a measurable function only of the
method-specific span, because its threshold is deterministic on this fiber.
Lemma~\ref{lem:step-001-haar-independence} supplies the product-Haar
conditional law.  Hence
\[
 \mathbb P(E_{\mathrm{cALS}}\cap E_{\mathrm{cGD}}\mid T=T_0)
 =\prod_{M\in\{\mathrm{cALS},\mathrm{cGD}\}}
 \mathbb P(E_M\mid T=T_0)
 \ge\frac12\cdot\frac12=\frac14.
\]
This includes the $T_0=0$ fiber, on which each factor is one.  The tower
property then gives
\[
 \begin{aligned}
 \mathbb P(E_{\mathrm{cALS}}\cap E_{\mathrm{cGD}})
 &=\mathbb E\!\left[
 \mathbb P(E_{\mathrm{cALS}}\cap E_{\mathrm{cGD}}\mid T)\right]\\
 &\ge\mathbb E\!\left[\frac14\right]=\frac14.
 \end{aligned}
\]
On the intersection, Lemma~\ref{lem:step-002-residual-margin} applies to
each method.  Only the two initial spans are conditionally independent; the
complete trajectories need not be independent.  The argument is pointwise in
$T_0$ and uses no concentration or nondegeneracy property of the tensor law.
The simultaneous rank-failure null set was already removed by
Lemma~\ref{lem:step-001-full-rank}.
\end{proof}

\subsection{Exact fixed-witness objective floor}\label{app:step-003}

\begin{proposition}[Exact same-target Pythagorean floor]\label{prop:step-003-pythagorean-floor}
Under Proposition~\ref{prop:step-001-fixed-span}, fix one of the two methods,
denote it by $M$, and fix an index $t$ for which its iterate is defined.  Set
\[
 R_M=(I-P_{\mathcal H_M})T.
\]
Then
\[
 T-S_t^M=R_M+(P_{\mathcal H_M}T-S_t^M),
 \qquad
 R_M\perp(P_{\mathcal H_M}T-S_t^M),
\]
and therefore
\[
 \|T-S_t^M\|_F^2
 =\|(I-P_{\mathcal H_M})T\|_F^2
 +\|P_{\mathcal H_M}T-S_t^M\|_F^2.
\]
Equivalently,
\[
 F_M(t)=\frac12\|R_M\|_F^2
 +\frac12\|P_{\mathcal H_M}T-S_t^M\|_F^2
 \ge\frac12\|(I-P_{\mathcal H_M})T\|_F^2.
\]
The bound holds at every defined index with the same fixed residual $R_M$.
\end{proposition}

\begin{proof}
Write $P=P_{\mathcal H_M}$.  Proposition~\ref{prop:step-001-fixed-span}
gives $PS_t^M=S_t^M$, so $PT-S_t^M\in\mathcal H_M$.  On the other hand,
\[
 PR_M=P(I-P)T=(P-P^2)T=0,
\]
and hence $R_M\in\mathcal H_M^\perp$.  Thus
$\langle R_M,PT-S_t^M\rangle_F=0$, while algebraically
\[
 T-S_t^M=(I-P)T+PT-S_t^M=R_M+(PT-S_t^M).
\]
Expanding the squared Frobenius norm proves the equality.  Multiplying by the
objective coefficient $1/2$ gives the exact objective decomposition, and
discarding only its displayed nonnegative in-subspace term gives the floor.
The residual $R_M$ depends only on $T$ and the fixed initialization subspace,
so no recurrence, accumulated error, horizon constant, or factor-size
condition is present.
\end{proof}

\begin{lemma}[Fixed normalized witness in the nonzero-residual branch]\label{lem:step-003-fixed-witness}
Under Proposition~\ref{prop:step-001-fixed-span} and
Proposition~\ref{prop:step-003-pythagorean-floor}, if $\|R_M\|_F>0$, then
\[
 W_M=\frac{R_M}{\|R_M\|_F}
\]
is a fixed unit tensor in $\mathcal H_M^\perp$ satisfying, for every defined
iterate,
\[
 \langle W_M,S_t^M\rangle_F=0,
 \qquad
 \langle W_M,T-S_t^M\rangle_F
 =\langle W_M,T\rangle_F=\|R_M\|_F,
\]
and hence $\|T-S_t^M\|_F\ge\|R_M\|_F$.  If $\|R_M\|_F=0$, no normalized
witness is defined or used, while
Proposition~\ref{prop:step-003-pythagorean-floor} remains valid.
\end{lemma}

\begin{proof}
In the positive-denominator branch, normalization gives
$\|W_M\|_F=1$ and, because $R_M\in\mathcal H_M^\perp$,
$W_M\in\mathcal H_M^\perp$.  Proposition~\ref{prop:step-001-fixed-span}
gives $S_t^M\in\mathcal H_M$, so
$\langle W_M,S_t^M\rangle_F=0$.  Also
\[
 \begin{aligned}
 \langle W_M,T\rangle_F
 &=\langle W_M,R_M+P_{\mathcal H_M}T\rangle_F\\
 &=\frac{\langle R_M,R_M\rangle_F}{\|R_M\|_F}
 =\|R_M\|_F.
 \end{aligned}
\]
Subtracting the zero inner product with $S_t^M$ proves the witness identity.
For completeness, Cauchy--Schwarz follows by applying nonnegativity of
$\|A-\lambda B\|_F^2$ at
$\lambda=\langle A,B\rangle_F/\|B\|_F^2$ when $B\ne0$; the zero case is
immediate.  Applying it with $A=W_M$ gives
\[
 \|R_M\|_F
 =|\langle W_M,T-S_t^M\rangle_F|
 \le\|T-S_t^M\|_F.
\]
Both $R_M$ and $W_M$ are fixed before the trajectory evolves.  When
$\|R_M\|_F=0$, the quotient is not assigned a value and the unconditional
Pythagorean identity is used instead.
\end{proof}

\begin{proposition}[Boundary-complete fixed obstruction]\label{prop:step-003-certificate}
Under Proposition~\ref{prop:step-001-fixed-span}, the Pythagorean floor in
Proposition~\ref{prop:step-003-pythagorean-floor} holds for every defined
iterate, including all zero-residual and degenerate factor cases, and the
normalized certificate in Lemma~\ref{lem:step-003-fixed-witness} is used only
on its positive-denominator branch.
\end{proposition}

\begin{proof}
The two cited results prove the asserted fixed floor and witness statement.
If $T=0$, then $R_M=0$ and the identity becomes
$\|S_t^M\|_F^2=0+\|S_t^M\|_F^2$.  If
$T\in\mathcal H_M$, the same zero-residual reduction applies.  The algebra
also remains valid at the literal full-span boundary $P_{\mathcal H_M}=I$,
although the current rank assumptions give $k<n$.  If
$S_t^M=P_{\mathcal H_M}T$, equality holds in the objective floor.  The
membership producer includes initialization and every first defined update,
so no burn-in is needed.  A nonexistent future GD iterate is not quantified
here; all-index existence is established later.  Finally, zero,
rank-deficient, or arbitrarily large finite factors do not change the
orthogonal decomposition once the represented tensor belongs to
$\mathcal H_M$.
\end{proof}

\subsection{Sequential constrained ALS descent and scalar convergence}\label{app:step-004}

\begin{lemma}[Exact CP mode-matricization identities]\label{lem:step-004-cp-matricization}
For every finite $X,Y,Z\in\mathbb R^{n\times k}$, under the fixed
Khatri--Rao ordering,
\[
 S(X,Y,Z)_{(1)}=X(Z\odot Y)^{\mathsf T},
 \quad
 S(X,Y,Z)_{(2)}=Y(Z\odot X)^{\mathsf T},
 \quad
 S(X,Y,Z)_{(3)}=Z(Y\odot X)^{\mathsf T}.
\]
Consequently, each matricized block least-squares objective is exactly the
tensor objective $F$.
\end{lemma}

\begin{proof}
The $i$th column of $Z\odot Y$ is $z_i\otimes y_i$, while the mode-1
matricization of $x_i\otimes y_i\otimes z_i$ is
$x_i(z_i\otimes y_i)^{\mathsf T}$.  Hence
\[
 \sum_{i=1}^k x_i(z_i\otimes y_i)^{\mathsf T}
 =X(Z\odot Y)^{\mathsf T}.
\]
The same coordinate calculation in modes 2 and 3 uses the ordered columns
$z_i\otimes x_i$ and $y_i\otimes x_i$.  Matricization permutes entries and
therefore preserves the Frobenius norm, proving exact equality of the matrix
and tensor residual objectives.
\end{proof}

\begin{lemma}[Singular-design Moore--Penrose least squares]\label{lem:step-004-pseudoinverse-ls}
Let $A\in\mathbb R^{p\times m}$ and $K\in\mathbb R^{m\times k}$ be finite.
Then
\[
 \phi(U)=\frac12\|A-UK^{\mathsf T}\|_F^2,
 \qquad U\in\mathbb R^{p\times k},
\]
is minimized by
\[
 U_\star=AK(K^{\mathsf T}K)^\dagger,
\]
and $U_\star$ is the unique minimizer of smallest Frobenius norm.  If
$K=U_K\Sigma V_K^{\mathsf T}$ is a compact SVD of positive rank and
$P_K=V_KV_K^{\mathsf T}$, every minimizer has the form
$U_\star+R(I_k-P_K)$.  If $K=0$, every $U$ minimizes $\phi$ and $U_\star=0$
is the unique minimum-norm minimizer.
\end{lemma}

\begin{proof}
First suppose $s=\operatorname{rank}(K)>0$, and extend $U_K,V_K$ to
orthogonal bases $[U_K\ U_0]$ of $\mathbb R^m$ and $[V_K\ V_0]$ of
$\mathbb R^k$.  Since $K^{\mathsf T}=V_K\Sigma U_K^{\mathsf T}$,
right multiplication by $[U_K\ U_0]$ gives
\[
 \begin{aligned}
 \|A-UK^{\mathsf T}\|_F^2
 &=\|(A-UV_K\Sigma U_K^{\mathsf T})[U_K\ U_0]\|_F^2\\
 &=\|AU_K-UV_K\Sigma\|_F^2+\|AU_0\|_F^2.
 \end{aligned}
\]
The second term is independent of $U$, and the first is minimized exactly
when
\[
 UV_K=AU_K\Sigma^{-1}.
\]
Thus all minimizers are
\[
 U=AU_K\Sigma^{-1}V_K^{\mathsf T}+R(I_k-P_K).
\]
The compact SVD also gives
\[
 \begin{aligned}
 K(K^{\mathsf T}K)^\dagger
 &=U_K\Sigma V_K^{\mathsf T}
 V_K\Sigma^{-2}V_K^{\mathsf T}\\
 &=U_K\Sigma^{-1}V_K^{\mathsf T},
 \end{aligned}
\]
so the first term is $U_\star$.  Orthogonality of the two right subspaces
gives
\[
 \|U\|_F^2=\|UV_K\|_F^2+\|U(I_k-P_K)\|_F^2.
\]
The unique smallest norm is obtained by setting the free component to zero.
If $s=0$, then $K=0$ and
$\phi(U)=\frac12\|A\|_F^2$ for every $U$, while the pseudoinverse formula
returns $U_\star=0$.  This proves the result at every rank without a
conditioning assumption.
\end{proof}

\begin{lemma}[Constrained orthogonal-basis block minimizer]\label{lem:step-004-constrained-block}
Let $A\in\mathbb R^{p\times m}$, $K\in\mathbb R^{m\times k}$, and
$Q\in\mathbb R^{p\times k}$ satisfy $Q^{\mathsf T}Q=I_k$.  Then
\[
 \min_{\operatorname{col}(X)\subseteq\operatorname{range}(Q)}
 \frac12\|A-XK^{\mathsf T}\|_F^2
\]
has the minimum-Frobenius-norm minimizer
\[
 X_\star=Q(Q^{\mathsf T}AK)(K^{\mathsf T}K)^\dagger.
\]
The conclusion holds for every rank of $K$, including $K=0$.
\end{lemma}

\begin{proof}
Every feasible matrix has a unique representation $X=QC$, with
$C=Q^{\mathsf T}X$.  Put $P=QQ^{\mathsf T}$.  Then
\[
 A-QCK^{\mathsf T}
 =(I_p-P)A+Q(Q^{\mathsf T}A-CK^{\mathsf T}).
\]
The two terms are Frobenius-orthogonal because $(I_p-P)Q=0$, and
$\|QB\|_F=\|B\|_F$.  Therefore
\[
 \|A-QCK^{\mathsf T}\|_F^2
 =\|(I_p-P)A\|_F^2
 +\|Q^{\mathsf T}A-CK^{\mathsf T}\|_F^2.
\]
The first term is independent of $C$.  Applying
Lemma~\ref{lem:step-004-pseudoinverse-ls} to $Q^{\mathsf T}A$ and $K$
gives
\[
 C_\star=(Q^{\mathsf T}A)K(K^{\mathsf T}K)^\dagger.
\]
Since $\|QC\|_F=\|C\|_F$, its minimum-norm property transfers to the
feasible matrix $X_\star$.  The same calculation includes the zero design.
\end{proof}

\begin{proposition}[Exact finite constrained ALS block updates]\label{prop:step-004-block-updates}
At any finite constrained ALS state with
$Q=Q_{\mathrm{cALS}}$, $Q^{\mathsf T}Q=I_k$, and
$X_t\in\operatorname{range}(Q)$, define in sequential order
\[
 K_t^x=Z_t\odot Y_t,
 \qquad K_t^y=Z_t\odot X_{t+1},
 \qquad K_t^z=Y_{t+1}\odot X_{t+1}.
\]
The three displayed updates are finite minimum-Frobenius-norm minimizers of
their current block problems, and
\[
 \begin{aligned}
 F(X_{t+1},Y_t,Z_t)&\le F(X_t,Y_t,Z_t),\\
 F(X_{t+1},Y_{t+1},Z_t)&\le F(X_{t+1},Y_t,Z_t),\\
 F(X_{t+1},Y_{t+1},Z_{t+1})&\le F(X_{t+1},Y_{t+1},Z_t).
 \end{aligned}
\]
\end{proposition}

\begin{proof}
By Lemma~\ref{lem:step-004-cp-matricization}, the $X$-block objective is
\[
 \frac12\|T_{(1)}-X(K_t^x)^{\mathsf T}\|_F^2.
\]
Lemma~\ref{lem:step-004-constrained-block}, with $A=T_{(1)}$ and
$K=K_t^x$, gives exactly
\[
 X_{t+1}=Q\left[Q^{\mathsf T}T_{(1)}K_t^x
 \bigl((K_t^x)^{\mathsf T}K_t^x\bigr)^\dagger\right].
\]
Because $X_t$ is feasible, the constrained minimum is no larger than its
value at $X_t$, proving the first inequality.

With $X_{t+1}$ fixed, Lemmas~\ref{lem:step-004-cp-matricization} and
\ref{lem:step-004-pseudoinverse-ls} give
\[
 Y_{t+1}=T_{(2)}K_t^y
 \bigl((K_t^y)^{\mathsf T}K_t^y\bigr)^\dagger
\]
as the minimum-norm $Y$-block minimizer, proving the second inequality.  With
$X_{t+1},Y_{t+1}$ fixed, the same lemmas give
\[
 Z_{t+1}=T_{(3)}K_t^z
 \bigl((K_t^z)^{\mathsf T}K_t^z\bigr)^\dagger
\]
as the minimum-norm $Z$-block minimizer, proving the third.  Every finite
Gram matrix has a finite Moore--Penrose inverse even when singular or zero,
so the selected updates are finite.
\end{proof}

\begin{proposition}[Well-defined sequential constrained ALS descent]\label{prop:step-004-sequential-descent}
For every finite constrained ALS initialization on the stipulated
orthonormal-basis domain, every sweep is defined,
$X_t\in\operatorname{range}(Q_{\mathrm{cALS}})$ for every $t$, and
\[
 F_{\mathrm{cALS}}(t+1)\le F_{\mathrm{cALS}}(t).
\]
More precisely, if
\[
 F_t^x=F(X_{t+1},Y_t,Z_t),
 \qquad
 F_t^y=F(X_{t+1},Y_{t+1},Z_t),
\]
then
\[
 F_{\mathrm{cALS}}(t+1)\le F_t^y\le F_t^x
 \le F_{\mathrm{cALS}}(t).
\]
\end{proposition}

\begin{proof}
At initialization,
\[
 X_0=G_x^{\mathrm{cALS}}
 \in\operatorname{range}(G_x^{\mathrm{cALS}})
 =\operatorname{range}(Q_{\mathrm{cALS}}).
\]
Assume $X_t,Y_t,Z_t$ are finite and $X_t$ is feasible.
Proposition~\ref{prop:step-004-block-updates} gives a finite
$X_{t+1}=Q_{\mathrm{cALS}}C_{t+1}$, so the new $X$ remains feasible.  Then
$K_t^y$ is finite and the same proposition gives finite $Y_{t+1}$; next
$K_t^z$ is finite and it gives finite $Z_{t+1}$.  This closes the induction.

The three block inequalities are composed in the actual update order:
\[
 \begin{aligned}
 F_{\mathrm{cALS}}(t+1)
 &=F(X_{t+1},Y_{t+1},Z_{t+1})\\
 &\le F_t^y\le F_t^x
 \le F(X_t,Y_t,Z_t)=F_{\mathrm{cALS}}(t).
 \end{aligned}
\]
The $Y$-block uses the newly minimized $X_{t+1}$, and the $Z$-block uses
both newly minimized factors.  No simultaneous update is substituted.
Nonuniqueness at a rank-deficient design does not change these inequalities,
because the selected pseudoinverse solution is still an exact minimizer.
\end{proof}

\begin{lemma}[Finite constrained ALS scalar limit and drop budget]\label{lem:step-004-scalar-limit}
Under Proposition~\ref{prop:step-004-sequential-descent},
\[
 L_{\mathrm{cALS}}
 :=\lim_{t\to\infty}F_{\mathrm{cALS}}(t)
 =\inf_{t\ge0}F_{\mathrm{cALS}}(t)
\]
exists as a finite real number.  For every $N\ge1$,
\[
 \sum_{t=0}^{N-1}
 \bigl(F_{\mathrm{cALS}}(t)-F_{\mathrm{cALS}}(t+1)\bigr)
 =F_{\mathrm{cALS}}(0)-F_{\mathrm{cALS}}(N)
 \le F_{\mathrm{cALS}}(0).
\]
\end{lemma}

\begin{proof}
The squared-loss definition gives $F_{\mathrm{cALS}}(t)\ge0$, and
Proposition~\ref{prop:step-004-sequential-descent} gives a nonincreasing
finite sequence.  Let
\[
 L=\inf_{t\ge0}F_{\mathrm{cALS}}(t),
 \qquad 0\le L\le F_{\mathrm{cALS}}(0)<\infty.
\]
For $\varepsilon>0$, choose $N$ with
$F_{\mathrm{cALS}}(N)<L+\varepsilon$.  For every $t\ge N$,
\[
 L\le F_{\mathrm{cALS}}(t)
 \le F_{\mathrm{cALS}}(N)<L+\varepsilon,
\]
which proves convergence to $L$.  Summing the successive differences
telescopes, and $F_{\mathrm{cALS}}(N)\ge0$ gives the stated budget.
No factor convergence or stationarity conclusion is used.
\end{proof}

\begin{proposition}[Constrained ALS scalar-limit certificate]\label{prop:step-004-certificate}
Under the finite constrained ALS definitions, every sweep is well-defined,
each displayed block update is an exact minimum-Frobenius-norm minimizer,
$F_{\mathrm{cALS}}(t)$ is nonnegative and nonincreasing, and its finite
scalar limit exists.  These conclusions require no Khatri--Rao rank,
conditioning, factor bound, or convergence-rate assumption.
\end{proposition}

\begin{proof}
Lemmas~\ref{lem:step-004-cp-matricization}--\ref{lem:step-004-constrained-block}
identify the actual block objectives and prove the singular-design formulas.
Propositions~\ref{prop:step-004-block-updates} and
\ref{prop:step-004-sequential-descent} give finite sequential descent, and
Lemma~\ref{lem:step-004-scalar-limit} gives the scalar limit.

If a design is zero, its block objective is independent of that factor and
the pseudoinverse selects the zero factor, the unique minimum-norm choice.  A
nonzero singular design retains only its row-space component.  The first
sweep is covered because $X_0$ is feasible by construction.  Tied minima and
zero drops are allowed.  If an objective value is zero, every later block
minimum remains zero by nonnegativity.  Factors or pseudoinverses may be
arbitrarily large at a finite time; the result asserts finite arrays at every
finite step but no uniform norm bound or continuity of the update map.
\end{proof}

\subsection{Coefficient Armijo descent and scalar convergence}\label{app:step-005}

Fix $Q=Q_{\mathrm{cGD}}$ and abbreviate
\[
 f_Q(C,Y,Z)=F(QC,Y,Z).
\]
We use the finite-dimensional product space
\[
 \mathcal U=\mathbb R^{k\times k}\times\mathbb R^{n\times k}
 \times\mathbb R^{n\times k}
\]
with inner product
\[
 \langle(C,Y,Z),(C',Y',Z')\rangle_{\mathcal U}
 =\langle C,C'\rangle_F+\langle Y,Y'\rangle_F+\langle Z,Z'\rangle_F.
\]
Its norm is exactly the block Frobenius norm used by the line search.

\begin{lemma}[Pointwise Taylor control for the coefficient loss]\label{lem:step-005-local-taylor}
If $u=(C,Y,Z)\in\mathcal U$ is finite and $g=\nabla f_Q(u)$, then $g$ is
finite and
\[
 L(u)=\max_{0\le s\le1}
 \|\nabla^2f_Q(u-sg)\|_{\mathrm{op}}<\infty.
\]
For every $0\le\eta\le1$,
\[
 f_Q(u-\eta g)
 \le f_Q(u)-\eta\|g\|_{\mathcal U}^2
 +\frac{L(u)\eta^2}{2}\|g\|_{\mathcal U}^2.
 \tag{A.1}\label{eq:step-005-taylor}
\]
\end{lemma}

\begin{proof}
For tensor indices $a,b,c$,
\[
 S(QC,Y,Z)_{abc}
 =\sum_{i=1}^k\sum_{\ell=1}^k
 Q_{a\ell}C_{\ell i}Y_{bi}Z_{ci}.
\]
Every entry of $T-S(QC,Y,Z)$ is therefore a polynomial in the entries of
$(C,Y,Z)$, and
\[
 f_Q(C,Y,Z)=\frac12\sum_{a,b,c}
 \bigl(T-S(QC,Y,Z)\bigr)_{abc}^2
\]
is a finite polynomial.  Thus it is $C^2$, and its gradient and Hessian have
finite entries at every finite state.

The map $s\mapsto u-sg$ sends the compact interval $[0,1]$ continuously to
a compact segment.  The function
$v\mapsto\|\nabla^2f_Q(v)\|_{\mathrm{op}}$ is continuous, so it attains a
finite maximum on that segment.  This proves the assertion about $L(u)$.

For $0\le\eta\le1$, define $\phi(s)=f_Q(u-sg)$.  Then
\[
 \phi'(0)=-\|g\|_{\mathcal U}^2,
 \qquad
 \phi''(s)=\langle g,\nabla^2f_Q(u-sg)g\rangle_{\mathcal U}.
\]
Applying the fundamental theorem of calculus twice gives the integral
remainder formula
\[
 \phi(\eta)=\phi(0)+\eta\phi'(0)
 +\int_0^\eta(\eta-s)\phi''(s)\,ds.
\]
Therefore
\[
 \begin{aligned}
 f_Q(u-\eta g)
 &=f_Q(u)-\eta\|g\|_{\mathcal U}^2
 +\int_0^\eta(\eta-s)
 \langle g,\nabla^2f_Q(u-sg)g\rangle_{\mathcal U}\,ds\\
 &\le f_Q(u)-\eta\|g\|_{\mathcal U}^2
 +\int_0^\eta(\eta-s)L(u)\|g\|_{\mathcal U}^2\,ds,
 \end{aligned}
\]
which is exactly \eqref{eq:step-005-taylor}.
\end{proof}

\begin{proposition}[Finite dyadic Armijo acceptance]\label{prop:step-005-armijo}
Under Lemma~\ref{lem:step-005-local-taylor}, if $u\in\mathcal U$ is finite,
$g=\nabla f_Q(u)$, and the line search tests $\eta_j=2^{-j}$, then an
acceptable trial occurs at a finite index.  The first acceptable $\eta$ is
well-defined, $u^+=u-\eta g$ is finite, and
\[
 f_Q(u^+)\le f_Q(u)-\frac{\eta}{2}\|g\|_{\mathcal U}^2.
 \tag{A.2}\label{eq:step-005-armijo}
\]
\end{proposition}

\begin{proof}
If $g=0$, the first trial $\eta_0=1$ leaves $u$ unchanged and satisfies the
Armijo inequality with equality.  Suppose $g\ne0$ and write $L=L(u)<\infty$.
If $L=0$, take $\eta_*=1$.  If $L>0$, set
\[
 j_*=\max\{0,\lceil\log_2L\rceil\},
 \qquad \eta_*=2^{-j_*}.
\]
Then $0<\eta_*\le1$ and $L\eta_*\le1$.  In both cases,
Lemma~\ref{lem:step-005-local-taylor} gives
\[
 \begin{aligned}
 f_Q(u-\eta_*g)
 &\le f_Q(u)-\eta_*\|g\|_{\mathcal U}^2
 +\frac{L\eta_*^2}{2}\|g\|_{\mathcal U}^2\\
 &=f_Q(u)-\eta_*
 \left(1-\frac{L\eta_*}{2}\right)\|g\|_{\mathcal U}^2\\
 &\le f_Q(u)-\frac{\eta_*}{2}\|g\|_{\mathcal U}^2.
 \end{aligned}
\]
Thus a trial no later than $j_*$ is acceptable.  The first acceptable trial
therefore occurs after finitely many tests and obeys
\eqref{eq:step-005-armijo}.  Since $u,g,$ and $\eta$ are finite, so is
$u^+=u-\eta g$.
\end{proof}

\begin{proposition}[Well-defined coefficient GD iterates]\label{prop:step-005-well-defined}
For the coefficient initialization above and the dyadic Armijo rule, every
integer $t\ge0$ has a finite state
\[
 u_t=(C_t,Y_t,Z_t),
\]
a finite accepted line-search index, and a step size $\eta_t\in(0,1]$ for
which
\[
 u_{t+1}=u_t-\eta_t\nabla f_Q(u_t)
\]
is finite and
\[
 f_Q(u_{t+1})
 \le f_Q(u_t)-\frac{\eta_t}{2}
 \|\nabla f_Q(u_t)\|_{\mathcal U}^2.
 \tag{A.3}\label{eq:step-005-alltime-armijo}
\]
Consequently $X_t=QC_t$ and the represented tensor are defined at every
finite index.
\end{proposition}

\begin{proof}
The formal initialization
\[
 u_0=(Q^{\mathsf T}G_x^{\mathrm{cGD}},
 G_y^{\mathrm{cGD}},G_z^{\mathrm{cGD}})
\]
is a finite element of $\mathcal U$ on the probability-one domain on which
the method is instantiated.  If $u_t$ is finite,
Proposition~\ref{prop:step-005-armijo} produces a finite accepted trial,
a finite $u_{t+1}$, and \eqref{eq:step-005-alltime-armijo}.  Induction proves
the assertion at every finite index.  Since $Q$ is fixed and finite,
$X_t=QC_t$ and all inputs to the next iteration are finite as well.
\end{proof}

\begin{proposition}[Monotone constrained GD objective and finite scalar limit]\label{prop:step-005-scalar-limit}
Under Proposition~\ref{prop:step-005-well-defined}, for every $t\ge0$,
\[
 0\le F_{\mathrm{cGD}}(t+1)\le F_{\mathrm{cGD}}(t)<\infty.
 \tag{A.4}\label{eq:step-005-monotone}
\]
For every $N\ge1$,
\[
 0\le\sum_{t=0}^{N-1}\frac{\eta_t}{2}
 \|\nabla f_Q(u_t)\|_{\mathcal U}^2
 \le F_{\mathrm{cGD}}(0)-F_{\mathrm{cGD}}(N)
 \le F_{\mathrm{cGD}}(0).
 \tag{A.5}\label{eq:step-005-budget}
\]
The finite scalar
\[
 L_{\mathrm{cGD}}=\inf_{t\ge0}F_{\mathrm{cGD}}(t)
 \in[0,F_{\mathrm{cGD}}(0)]
\]
satisfies
\[
 \lim_{t\to\infty}F_{\mathrm{cGD}}(t)=L_{\mathrm{cGD}}.
 \tag{A.6}\label{eq:step-005-limit}
\]
\end{proposition}

\begin{proof}
The coefficient and actual objectives agree exactly:
\[
 f_Q(u_t)=F(QC_t,Y_t,Z_t)=F_{\mathrm{cGD}}(t).
 \tag{A.7}\label{eq:step-005-objective-identity}
\]
This finite polynomial is one half of a squared Frobenius norm and is
therefore nonnegative.  Substituting
\eqref{eq:step-005-objective-identity} into
\eqref{eq:step-005-alltime-armijo} proves
\eqref{eq:step-005-monotone}.  Summation telescopes and gives
\eqref{eq:step-005-budget}.

The nonincreasing sequence is bounded below by zero.  For
$L_{\mathrm{cGD}}=\inf_tF_{\mathrm{cGD}}(t)$ and $\varepsilon>0$, choose
$N$ with
$F_{\mathrm{cGD}}(N)<L_{\mathrm{cGD}}+\varepsilon$.  For $t\ge N$,
\[
 L_{\mathrm{cGD}}\le F_{\mathrm{cGD}}(t)
 \le F_{\mathrm{cGD}}(N)<L_{\mathrm{cGD}}+\varepsilon,
\]
which proves \eqref{eq:step-005-limit}.

If $\nabla f_Q(u_t)=0$, the first trial is accepted and the state is
stationary.  If $f_Q(u_t)=0$, differentiability and nonnegativity force
$\nabla f_Q(u_t)=0$: otherwise, as $s\downarrow0$,
\[
 f_Q(u_t-s\nabla f_Q(u_t))
 =-s\|\nabla f_Q(u_t)\|_{\mathcal U}^2+o(s)<0,
\]
a contradiction.  Thus an exact zero objective remains zero.
\end{proof}

\begin{proposition}[Coefficient Armijo scalar-limit certificate]\label{prop:step-005-certificate}
For constrained GD in coefficient variables, every dyadic line search
terminates and every finite-index iterate exists.  Its objective values are
nonnegative and nonincreasing, and their finite scalar limit exists.  No
trajectory boundedness, global Lipschitz constant, uniform Hessian bound,
positive all-time step-size lower bound, or parameter convergence is required.
\end{proposition}

\begin{proof}
Lemma~\ref{lem:step-005-local-taylor} supplies a finite Hessian bound only on
the current search segment.  Proposition~\ref{prop:step-005-armijo} converts
that bound into finite one-step acceptance, and
Proposition~\ref{prop:step-005-well-defined} closes the local finiteness
condition by induction.  Proposition~\ref{prop:step-005-scalar-limit}
transfers the accepted coefficient decrease to the actual tensor objective,
telescopes the decreases, and proves scalar convergence.

The local numbers $L(u_t)$ may grow without bound, Armijo step sizes may
tend to zero, and factor norms may be unbounded over time.  None of these
possibilities prevents finite acceptance at a fixed finite state.  The same
proof applies at $T=0$, at zero gradient, and at zero objective, with exact
preservation of the zero baseline.
\end{proof}

\subsection{Why the fixed witness does not transfer to unconstrained updates}\label{app:step-006}

\begin{proposition}[Exact first-update leakage for unconstrained GD and ALS]\label{prop:step-006-first-update}
On the probability-one full-rank domain of
Lemma~\ref{lem:step-001-full-rank}, fix either method-specific initial triple
and suppress its method superscript.  Let
\[
 \mathcal S_0=\operatorname{range}(X_0),
 \qquad P_0=P_{\mathcal S_0},
 \qquad K_0=Z_0\odot Y_0.
\]
For one full-variable Euclidean gradient step with finite step size $\eta_0$,
\[
 X_1^{\mathrm{uGD}}
 =X_0-\eta_0\nabla_XF(X_0,Y_0,Z_0),
\]
one has
\[
 (I-P_0)X_1^{\mathrm{uGD}}
 =\eta_0(I-P_0)T_{(1)}K_0.
 \tag{A.8}\label{eq:step-006-gd-leakage}
\]
For the minimum-Frobenius-norm unconstrained ALS $X$-block update,
\[
 X_1^{\mathrm{uALS}}
 =T_{(1)}K_0(K_0^{\mathsf T}K_0)^\dagger,
\]
one has
\[
 (I-P_0)X_1^{\mathrm{uALS}}
 =(I-P_0)T_{(1)}K_0(K_0^{\mathsf T}K_0)^\dagger.
 \tag{A.9}\label{eq:step-006-als-leakage}
\]
Thus first-step fixed-span preservation is equivalent to the vanishing of the
corresponding displayed forcing term; it is not implied merely by
$X_0\in\mathcal S_0$.
\end{proposition}

\begin{proof}
Lemma~\ref{lem:step-004-cp-matricization} gives
\[
 S(X,Y_0,Z_0)_{(1)}=XK_0^{\mathsf T},
 \qquad
 F(X,Y_0,Z_0)=\frac12\|T_{(1)}-XK_0^{\mathsf T}\|_F^2.
\]
For any $H\in\mathbb R^{n\times k}$,
\[
 \begin{aligned}
 \frac{d}{d\epsilon}F(X+\epsilon H,Y_0,Z_0)\bigg|_{\epsilon=0}
 &=\langle XK_0^{\mathsf T}-T_{(1)},HK_0^{\mathsf T}\rangle_F\\
 &=\langle(XK_0^{\mathsf T}-T_{(1)})K_0,H\rangle_F.
 \end{aligned}
\]
Therefore
\[
 \nabla_XF(X_0,Y_0,Z_0)
 =X_0K_0^{\mathsf T}K_0-T_{(1)}K_0,
\]
and
\[
 X_1^{\mathrm{uGD}}
 =X_0(I_k-\eta_0K_0^{\mathsf T}K_0)+\eta_0T_{(1)}K_0.
\]
Since $(I-P_0)X_0=0$, left multiplication by $I-P_0$ proves
\eqref{eq:step-006-gd-leakage}.  The calculation is unchanged if the other
two factor components are updated simultaneously, because it uses the
$X$-gradient at the common initial triple.

For ALS, let $K_0=U\Sigma V^{\mathsf T}$ be an SVD.  The same rowwise
least-squares calculation as in
Lemma~\ref{lem:step-004-pseudoinverse-ls} gives the minimum-norm minimizer
$T_{(1)}(K_0^{\mathsf T})^\dagger$.  The SVD yields
\[
 (K_0^{\mathsf T})^\dagger
 =U\Sigma^\dagger V^{\mathsf T}
 =K_0(K_0^{\mathsf T}K_0)^\dagger,
\]
including rank-deficient and zero $K_0$.  Applying $I-P_0$ proves
\eqref{eq:step-006-als-leakage}.

The identities do not assert that either right-hand side is always nonzero.
It may vanish when $K_0=0$, when the target contraction lies in
$\mathcal S_0$, through rank deficiency, or by accidental cancellation.
Conversely, neither unconstrained update contains the fixed left factor that
enforces zero leakage in Proposition~\ref{prop:step-001-fixed-span}.
\end{proof}

\begin{proposition}[Exact outside-span residual decomposition]\label{prop:step-006-residual-decomposition}
Let $P_0$ be any orthogonal projector on $\mathbb R^n$, and define
\[
 \mathcal H_0=\operatorname{range}(P_0)\otimes\mathbb R^n\otimes\mathbb R^n,
 \qquad
 \Pi_0=P_0\otimes I_n\otimes I_n.
\]
For arbitrary finite factors $X_t,Y_t,Z_t$, set
\[
 S_t=S(X_t,Y_t,Z_t),
 \qquad
 R_0=(I-\Pi_0)T,
 \qquad
 \Lambda_t=(I-\Pi_0)S_t.
\]
Then
\[
 \Lambda_t=S((I-P_0)X_t,Y_t,Z_t),
 \tag{A.10}\label{eq:step-006-tensor-leakage}
\]
and
\[
 T-S_t=(R_0-\Lambda_t)+(\Pi_0T-\Pi_0S_t),
 \tag{A.11}\label{eq:step-006-residual-split}
\]
is an orthogonal decomposition.  Consequently,
\[
 \|T-S_t\|_F^2
 =\|R_0-\Lambda_t\|_F^2
 +\|\Pi_0T-\Pi_0S_t\|_F^2,
 \tag{A.12}\label{eq:step-006-residual-pythagoras}
\]
and
\[
 F(X_t,Y_t,Z_t)
 =\frac12\|R_0-\Lambda_t\|_F^2
 +\frac12\|\Pi_0T-\Pi_0S_t\|_F^2.
 \tag{A.13}\label{eq:step-006-objective-split}
\]
If $\|R_0\|_F>0$ and $W_0=R_0/\|R_0\|_F$, then
\[
 \langle W_0,T-S_t\rangle_F
 =\|R_0\|_F-\langle W_0,\Lambda_t\rangle_F.
 \tag{A.14}\label{eq:step-006-witness-subtraction}
\]
\end{proposition}

\begin{proof}
Writing factor columns as $x_{t,i},y_{t,i},z_{t,i}$ and using
$I-\Pi_0=(I-P_0)\otimes I_n\otimes I_n$ gives
\[
 \begin{aligned}
 (I-\Pi_0)S_t
 &=\sum_{i=1}^k((I-P_0)x_{t,i})\otimes y_{t,i}\otimes z_{t,i}\\
 &=S((I-P_0)X_t,Y_t,Z_t),
 \end{aligned}
\]
which proves \eqref{eq:step-006-tensor-leakage}.  Nonzero leaking factor
columns may cancel in the represented tensor, so no converse is asserted.

Insert $I=\Pi_0+(I-\Pi_0)$ separately on $T$ and $S_t$:
\[
 \begin{aligned}
 T-S_t
 &=(I-\Pi_0)T-(I-\Pi_0)S_t+\Pi_0T-\Pi_0S_t\\
 &=(R_0-\Lambda_t)+(\Pi_0T-\Pi_0S_t).
 \end{aligned}
\]
The first summand lies in $\mathcal H_0^\perp$ and the second in
$\mathcal H_0$.  Their inner product vanishes, proving
\eqref{eq:step-006-residual-pythagoras} and, after multiplication by $1/2$,
\eqref{eq:step-006-objective-split}.

In the branch $\|R_0\|_F>0$, $W_0$ is a unit tensor in
$\mathcal H_0^\perp$ and is orthogonal to the in-span residual.  Therefore
\[
 \begin{aligned}
 \langle W_0,T-S_t\rangle_F
 &=\langle W_0,R_0-\Lambda_t\rangle_F\\
 &=\|R_0\|_F-\langle W_0,\Lambda_t\rangle_F.
 \end{aligned}
\]
Orthogonal projection supplies no sign for the last inner product.
\end{proof}

\begin{proposition}[Fixed-witness non-transfer certificate]\label{prop:step-006-nontransfer}
Under Propositions~\ref{prop:step-001-fixed-span},
\ref{prop:step-003-pythagorean-floor},
\ref{prop:step-006-first-update}, and
\ref{prop:step-006-residual-decomposition}, the fixed-witness proof for the
two constrained methods does not transfer to the ordinary unconstrained
comparison dynamics from these interfaces alone.  More precisely:
\begin{enumerate}[label=(\roman*)]
 \item For the constrained methods, $\Lambda_t=0$, so the fixed term
 $\|R_0\|_F^2$ is recovered exactly.
 \item For either unconstrained comparison,
 \[
 F(X_t,Y_t,Z_t)
 \ge\frac12\|R_0-\Lambda_t\|_F^2
 \ge\frac12\bigl(\max\{\|R_0\|_F-\|\Lambda_t\|_F,0\}\bigr)^2.
 \]
 \item If an additional estimate
 $\|\Lambda_t\|_F\le\gamma\|R_0\|_F$ held for some $0\le\gamma<1$, then
 \[
 F(X_t,Y_t,Z_t)\ge\frac12(1-\gamma)^2\|R_0\|_F^2.
 \]
 No such estimate follows from the preceding results for unconstrained
 iterates.
 \item If $R_0-\Lambda_t\ne0$, the endogenous tensor
 \[
 \widetilde W_t=\frac{R_0-\Lambda_t}{\|R_0-\Lambda_t\|_F}
 \in\mathcal H_0^\perp
 \]
 satisfies
 \[
 \langle\widetilde W_t,T-S_t\rangle_F=\|R_0-\Lambda_t\|_F.
 \]
 It yields a positive theorem-level floor only after a separate lower bound
 on its generated residual norm.
\end{enumerate}
These statements show only that the constrained proof interface is absent.
They do not assert that unconstrained iterates always leak, that cancellation
occurs, or that unconstrained ALS or gradient descent succeeds or fails.
\end{proposition}

\begin{proof}
Proposition~\ref{prop:step-001-fixed-span} gives $\Lambda_t=0$ in the
constrained methods.  Substitution in
Proposition~\ref{prop:step-006-residual-decomposition} reproduces the exact
Pythagorean floor of Proposition~\ref{prop:step-003-pythagorean-floor}.

For an unconstrained represented tensor, discard only the nonnegative
in-span term in \eqref{eq:step-006-objective-split}.  The reverse triangle
inequality gives
\[
 \|R_0-\Lambda_t\|_F
 \ge\max\{\|R_0\|_F-\|\Lambda_t\|_F,0\},
\]
which proves the second assertion.  Substituting the displayed conditional
$\gamma$-bound proves the third.  This parameter identifies one sufficient
missing interface; it is not an assumption of the constrained theorem.

When $R_0-\Lambda_t\ne0$, both terms lie in $\mathcal H_0^\perp$, so the
normalized difference lies there and is orthogonal to the in-span residual.
Pairing it with \eqref{eq:step-006-residual-split} proves the endogenous
witness identity.  Its right-hand side can vanish or be arbitrarily smaller
than $\|R_0\|_F$ without additional trajectory control.

Equations~\eqref{eq:step-006-gd-leakage} and
\eqref{eq:step-006-als-leakage} show that zero leakage is not enforced by an
unconstrained update; it requires an additional equality involving the target
contraction.  Accidental first-step vanishing is not an all-time invariant.

All degenerate regimes are consistent with this conclusion.  If $T=0$, then
$R_0=0$, both displayed first-update forcing terms vanish, and no normalized
fixed witness is defined.  If $K_0=0$, both first-update leakage terms also
vanish; singular $K_0$ is covered by the pseudoinverse.  Factor leakage need
not imply tensor leakage because rank-one terms can vanish or cancel.  Exact
complement cancellation $R_0=\Lambda_t$ is algebraically allowed but not
claimed to occur.  If $R_0=0$ or $P_0=I_n$, every outside-span expression
reduces to zero and normalization is omitted.  The endogenous witness is
likewise omitted at its zero denominator.  No smoothing-error term appears:
the non-transfer issue is exactly the absence of an unconstrained span
control in the same Frobenius geometry.
\end{proof}

\subsection{Joint probability and positive-limit assembly}\label{app:step-007}

\begin{proposition}[Joint-event residual and all-time objective floors]\label{prop:step-007-joint-floor}
Under Assumptions~\ref{assump:dimension} and \ref{assump:rank_window},
Proposition~\ref{prop:step-002-joint-event},
Proposition~\ref{prop:step-003-pythagorean-floor},
Proposition~\ref{prop:step-004-sequential-descent}, and
Proposition~\ref{prop:step-005-well-defined}, define
\[
 \mathcal E=E_{\mathrm{cALS}}\cap E_{\mathrm{cGD}}.
\]
On $\mathcal E$, for every $M\in\{\mathrm{cALS},\mathrm{cGD}\}$ and every
integer $t\ge0$,
\[
 \|(I-P_{\mathcal H_M})T\|_F^2
 \ge\left(1-2\frac{k}{n}\right)\|T\|_F^2
 \ge\frac34\|T\|_F^2
 \tag{A.15}\label{eq:step-007-residual-floor}
\]
and
\[
 F_M(t)\ge\frac38\|T\|_F^2.
 \tag{A.16}\label{eq:step-007-objective-floor}
\]
If $T\ne0$, then $W_M$ is well-defined, fixed in $t$, and certifies the same
residual obstruction at every finite index.  If $T=0$, no normalized witness
is used and both inequalities remain valid with zero right-hand side.
\end{proposition}

\begin{proof}
The dimension and rank assumptions give
\[
 \frac{k}{n}\le\frac{r^{5/4}}{8r^{5/4}}=\frac18.
 \tag{A.17}\label{eq:step-007-ratio}
\]
Proposition~\ref{prop:step-002-joint-event} gives on $\mathcal E$, for both
methods,
\[
 \|(I-P_{\mathcal H_M})T\|_F^2
 \ge\left(1-2\frac{k}{n}\right)\|T\|_F^2.
\]
Equation~\eqref{eq:step-007-ratio} gives
$1-2k/n\ge1-2(1/8)=3/4$, proving
\eqref{eq:step-007-residual-floor}, including equality at the boundary.

Proposition~\ref{prop:step-004-sequential-descent} proves that every
constrained ALS index exists, and
Proposition~\ref{prop:step-005-well-defined} proves the same for constrained
GD.  Therefore the defined-index scope of
Proposition~\ref{prop:step-003-pythagorean-floor} covers every integer
$t\ge0$.  Combining its exact floor with
\eqref{eq:step-007-residual-floor} yields
\[
 \begin{aligned}
 F_M(t)
 &\ge\frac12\|(I-P_{\mathcal H_M})T\|_F^2\\
 &\ge\frac12\cdot\frac34\|T\|_F^2
 =\frac38\|T\|_F^2.
 \end{aligned}
\]
The only discarded quantity is the explicit nonnegative in-subspace squared
residual in the exact Pythagorean identity.

If $T\ne0$, then
\[
 \|(I-P_{\mathcal H_M})T\|_F^2
 \ge\frac34\|T\|_F^2>0,
\]
so Lemma~\ref{lem:step-003-fixed-witness} applies and gives
\[
 \langle W_M,T-S_t^M\rangle_F
 =\|(I-P_{\mathcal H_M})T\|_F
 \quad\text{for every }t\ge0.
\]
At $T=0$ the quotient is omitted and the unnormalized identity gives the
zero floor.
\end{proof}

\begin{proposition}[Passage to both objective limits]\label{prop:step-007-limit-passage}
Under Proposition~\ref{prop:step-007-joint-floor},
Lemma~\ref{lem:step-004-scalar-limit}, and
Proposition~\ref{prop:step-005-scalar-limit}, the finite limits
\[
 \ell_{\mathrm{cALS}}=\lim_{t\to\infty}F_{\mathrm{cALS}}(t),
 \qquad
 \ell_{\mathrm{cGD}}=\lim_{t\to\infty}F_{\mathrm{cGD}}(t)
\]
exist, and on $\mathcal E$ satisfy simultaneously
\[
 \ell_M\ge\frac38\|T\|_F^2,
 \qquad M\in\{\mathrm{cALS},\mathrm{cGD}\}.
 \tag{A.18}\label{eq:step-007-limit-floor}
\]
\end{proposition}

\begin{proof}
Lemma~\ref{lem:step-004-scalar-limit} gives existence and finiteness of the
constrained ALS limit, and
Proposition~\ref{prop:step-005-scalar-limit} gives the constrained GD limit.
On $\mathcal E$, Proposition~\ref{prop:step-007-joint-floor} gives
\[
 F_M(t)\ge c_T:=\frac38\|T\|_F^2
 \quad\text{for every }t\ge0.
\]
If $\ell_M<c_T$, choose
$\varepsilon=(c_T-\ell_M)/2>0$.  Convergence would give an $N$ such that
\[
 F_M(t)<\ell_M+\varepsilon
 =\frac{\ell_M+c_T}{2}<c_T
 \quad\text{for every }t\ge N,
\]
contradicting the all-time floor.  Thus
\eqref{eq:step-007-limit-floor} holds for each method.  No limit is
interchanged with expectation, probability, or optimization.
\end{proof}

\begin{proposition}[Quantitative specialization and constrained positive-limit theorem]\label{prop:step-007-material-partial}
Fix $q>0$ and set
\[
 \alpha=\frac14,
 \qquad
 L(r)=r^{1+\alpha}=r^{5/4},
 \qquad
 r_0=1.
\]
Under Assumptions~\ref{assump:dimension}--\ref{assump:joint_initialization},
for every integer $r\ge r_0$, every integer $n\ge8L(r)$, every integer
$r<k\le L(r)$, and every unrestricted deterministic
$\bar A,\bar B,\bar C\in\mathbb R^{n\times r}$, draw
\[
 a_j=\bar a_j+\xi_j^a,
 \qquad b_j=\bar b_j+\xi_j^b,
 \qquad c_j=\bar c_j+\xi_j^c,
\]
where all perturbations are independent and
\[
 \xi_j^a,\xi_j^b,\xi_j^c
 \sim\mathcal N\!\left(0,\frac{r^{-2q}}{n}I_n\right),
\]
and set $T=\sum_{j=1}^r a_j\otimes b_j\otimes c_j$.  Independently draw
the two complete Gaussian initialization triples as in
Assumption~\ref{assump:joint_initialization}, use the same $T$ for both
methods, and run exactly the constrained procedures in the theoretical setup.
Then
\[
 \mathbb P\!\left[
 \bigcap_{M\in\{\mathrm{cALS},\mathrm{cGD}\}}
 \left\{
 \lim_{t\to\infty}F_M(t)\text{ exists as a finite real number and }
 \lim_{t\to\infty}F_M(t)\ge\frac38\|T\|_F^2
 \right\}
 \right]\ge\frac14.
 \tag{A.19}\label{eq:step-007-final-event}
\]
On the underlying probability-$1/4$ event, the residual-energy and all-time
objective bounds \eqref{eq:step-007-residual-floor}--
\eqref{eq:step-007-objective-floor} also hold for both methods.  The constants
hide no dependence on $r,n,k,q,\rho,T$, the deterministic bases, either
initialization, the method, the iteration, factor norms, design conditioning,
or line-search history.  The conclusion is pointwise uniform over each fixed
deterministic base triple and concerns only the displayed constrained
procedures; it makes no assertion for ordinary unconstrained ALS or
full-variable gradient descent.
\end{proposition}

\begin{proof}
Fix admissible $q,r,n,k$ and a deterministic base triple, and fix an arbitrary
possible tensor value $T_0$.  Proposition~\ref{prop:step-002-joint-event}
gives under the regular conditional initialization law
\[
 \mathbb P(\mathcal E\mid T=T_0)\ge\frac14.
 \tag{A.20}\label{eq:step-007-conditional-event}
\]
The probability-one Gaussian full-rank domain can be intersected with
$\mathcal E$ without changing this bound.  The constrained ALS and GD
existence results are pathwise and introduce no further random event.

On $\mathcal E$, Proposition~\ref{prop:step-007-limit-passage} gives both
limit conclusions in \eqref{eq:step-007-final-event}.  Therefore, for every
$T_0$,
\[
 \begin{aligned}
 &\mathbb P\!\left[
 \bigcap_{M\in\{\mathrm{cALS},\mathrm{cGD}\}}
 \left\{
 \lim_{t\to\infty}F_M(t)\text{ exists and }
 \lim_{t\to\infty}F_M(t)\ge\frac38\|T\|_F^2
 \right\}
 \middle|T=T_0\right]\\
 &\qquad\ge\mathbb P(\mathcal E\mid T=T_0)\ge\frac14.
 \end{aligned}
 \tag{A.21}\label{eq:step-007-conditional-success}
\]
This is pointwise in $T_0$ and does not use a favorable smoothing event or
independence of the complete trajectories.

Averaging \eqref{eq:step-007-conditional-success} over the smoothing law and
using the tower property gives
\[
 \begin{aligned}
 &\mathbb P\!\left[
 \bigcap_{M\in\{\mathrm{cALS},\mathrm{cGD}\}}
 \left\{
 \lim_{t\to\infty}F_M(t)\text{ exists and }
 \lim_{t\to\infty}F_M(t)\ge\frac38\|T\|_F^2
 \right\}\right]\\
 &\quad=\mathbb E_T\!\left[
 \mathbb P\!\left(
 \bigcap_M\left\{
 \lim_{t\to\infty}F_M(t)\text{ exists and }
 \lim_{t\to\infty}F_M(t)\ge\frac38\|T\|_F^2
 \right\}\middle|T\right)\right]\\
 &\quad\ge\mathbb E_T\!\left[\frac14\right]=\frac14.
 \end{aligned}
\]
Because the conditional argument holds for every $T_0$, no property of the
deterministic bases and no smoothing concentration, conditioning, or
incoherence event is used.  The same numerical constants apply after fixing
each deterministic base triple; no simultaneous event over all triples is
claimed.

The complete quantitative specialization is explicit:
\[
 k\le L(r)=r^{5/4},
 \qquad n\ge8L(r),
 \qquad \frac{k}{n}\le\frac18,
\]
and therefore
\[
 1-2\frac{k}{n}\ge\frac34,
 \qquad
 \frac12\left(1-2\frac{k}{n}\right)\ge\frac38.
\]
The two conditional method probabilities satisfy
$\frac12\cdot\frac12=\frac14$, and the tower conversion preserves this
constant.  No term is hidden or absorbed.  The arithmetic is valid for every
positive integer $r$, so $r_0=1$ is admissible; if the rank window is empty
for a small $r$, the universal assertion over admissible $k$ is vacuous.
The identity $L(r)=r^{1+\alpha}$ gives $\alpha=1/4$.

At $T=0$, both projection events are sure, both scalar limits exist, and the
required lower bound is exactly zero; no normalized witness is needed.  For
nonzero $T$ on $\mathcal E$, both witness denominators are positive.  At
$k/n=1/8$, the constants are exactly $3/4$ and $3/8$.  Zero block drops,
singular constrained ALS designs, and zero constrained GD gradients are
already included in the scalar-limit results.  The shared target prevents an
unconditional independence assertion, but conditional independence and the
tower property suffice.

Finally, Proposition~\ref{prop:step-006-nontransfer} shows that ordinary
unconstrained updates introduce the generated complement component
$\Lambda_t$ and do not inherit the fixed-witness interface.  Thus the theorem
neither states nor implies an unconstrained positive-limit result.
\end{proof}

\subsection{Proof of the Main Theorem}

\begin{proof}[Proof of Theorem~\ref{thm:main}]
Proposition~\ref{prop:step-007-material-partial} verifies the choices
$r_0=1$, $\alpha=1/4$, and $L(r)=r^{5/4}$; all dimension and rank
conditions; the exact $3/4$ residual-energy and $3/8$ objective constants;
the conditional product probability and tower conversion giving confidence
$1/4$; the existence of both finite scalar limits; and pointwise uniformity
over unrestricted deterministic bases.  Its event also contains the all-time
floors from Proposition~\ref{prop:step-007-joint-floor}.  These conclusions
are precisely those asserted in Theorem~\ref{thm:main}.  The constrained-only
scope follows from Proposition~\ref{prop:step-006-nontransfer}, which exposes
the uncontrolled outside-span term for ordinary unconstrained updates without
making an asymptotic claim about those dynamics.
\end{proof}
